% Preâmbulo comum aos documentos do Grupo VI (classe abnTeX2, normas ABNT).
% Compilar com Tectonic (XeTeX): ver gerar_pdf.sh.
\documentclass[12pt,openright,oneside,a4paper,chapter=TITLE,english,brazil]{abntex2}

% ---------------------------------------------------------------- Fontes
\usepackage{fontspec}
\setmainfont{texgyretermes}[
  Extension=.otf,
  UprightFont=*-regular, BoldFont=*-bold,
  ItalicFont=*-italic, BoldItalicFont=*-bolditalic]
\setmonofont{DejaVuSansMono}[
  Extension=.ttf, BoldFont=*-Bold, Scale=0.82]
\setsansfont{texgyreheros}[
  Extension=.otf,
  UprightFont=*-regular, BoldFont=*-bold,
  ItalicFont=*-italic, BoldItalicFont=*-bolditalic]
\usepackage{amsmath}

% ---------------------------------------------------------------- Pacotes
\usepackage{indentfirst}
\usepackage{graphicx}
\usepackage{xcolor}
\usepackage{booktabs}
\usepackage{array}
\usepackage{tabularx}
\usepackage{longtable}
\usepackage{xltabular}
\usepackage{multirow}
\usepackage[section]{placeins}
\usepackage{enumitem}
\usepackage{listings}
\usepackage{tikz}
\usepackage{adjustbox}
\usetikzlibrary{arrows.meta,positioning,fit,calc,shapes.geometric,backgrounds}
\usepackage[alf,abnt-emphasize=bf]{abntex2cite}

\graphicspath{{../../backend/resultados/}}

% ---------------------------------------------------------------- Aparência
\definecolor{cinzacodigo}{gray}{0.96}
\definecolor{cinzaborda}{gray}{0.75}
\lstset{
  language=Python,
  basicstyle=\ttfamily\small,
  keywordstyle=\bfseries,
  commentstyle=\itshape\color{black!65},
  stringstyle=\color{black!80},
  backgroundcolor=\color{cinzacodigo},
  frame=single, rulecolor=\color{cinzaborda},
  framesep=4pt, xleftmargin=6pt, xrightmargin=6pt,
  columns=fullflexible, keepspaces=true,
  showstringspaces=false, breaklines=true,
  numbers=none, aboveskip=0.8em, belowskip=0.8em,
}
\lstdefinestyle{json}{language={}, basicstyle=\ttfamily\footnotesize, keywordstyle={}}
\renewcommand{\lstlistingname}{Código}
\AtBeginDocument{\renewcommand{\thelstlisting}{\arabic{lstlisting}}}
\renewcommand{\lstlistlistingname}{Lista de códigos}

\setlist{itemsep=0.2em, topsep=0.3em}
\setlist[itemize]{label=--}
\newcolumntype{L}[1]{>{\raggedright\arraybackslash}p{#1}}
\newcolumntype{Y}{>{\raggedright\arraybackslash}X}
\newcolumntype{C}[1]{>{\centering\arraybackslash}p{#1}}
\renewcommand{\arraystretch}{1.25}
\setlength{\arrayrulewidth}{0.5pt}
% Quadros (ABNT): texto, bordas fechadas. Tabelas (IBGE): dados numéricos, laterais abertas.
% Definição do ambiente quadro conforme a documentação do abnTeX2.
\newcommand{\quadroname}{Quadro}
\newcommand{\listofquadrosname}{Lista de quadros}
\newcommand{\quadroautorefname}{Quadro}
\newfloat[chapter]{quadro}{loq}{\quadroname}
\newlistof{listofquadros}{loq}{\listofquadrosname}
\newlistentry{quadro}{loq}{0}
\setfloatadjustment{quadro}{\centering}
\counterwithout{quadro}{chapter}
\renewcommand{\cftquadroname}{\quadroname\space}
\renewcommand*{\cftquadroaftersnum}{\hfill--\hfill}
\setfloatlocations{quadro}{hbtp}

% Fonte de figuras e tabelas produzidas pelo grupo
\newcommand{\fonteautores}{\fonte{Elaborado pelos autores (2026).}}
% Diagrama TikZ: \diagrama{arquivo} reduz a figura à largura do texto, se preciso
\newcommand{\diagrama}[1]{\begin{adjustbox}{max width=\textwidth}\input{figuras/#1}\end{adjustbox}\par}
% Nome de arquivo, função ou campo do código
% (permite quebra de linha depois de \_ e de /)
\newcommand{\codsublinhado}{\textunderscore\allowbreak}
\newcommand{\cod}[1]{{\ttfamily\let\_\codsublinhado#1}}

% ---------------------------------------------------------------- PDF
\makeatletter
\hypersetup{
  pdftitle={\@title},
  pdfauthor={Grupo VI -- Sistemas Distribuídos -- UNIFESSPA},
  pdfsubject={Agregador de Preços de Eletrônicos},
  pdfcreator={LaTeX com abnTeX2},
  colorlinks=false, hidelinks, bookmarksdepth=4}
\makeatother

\setlength{\parindent}{1.25cm}
\setlength{\parskip}{0.2cm}
\setlength{\emergencystretch}{3em}

% Títulos em Times, negrito (seções primárias em caixa alta, como pede a NBR 6024)
\renewcommand{\ABNTEXchapterfont}{\rmfamily\bfseries}
\renewcommand{\ABNTEXchapterfontsize}{\normalsize}
\renewcommand{\ABNTEXsectionfont}{\rmfamily\bfseries}
\renewcommand{\ABNTEXsectionfontsize}{\normalsize}
\renewcommand{\ABNTEXsubsectionfont}{\rmfamily\bfseries}
\renewcommand{\ABNTEXsubsectionfontsize}{\normalsize}
\renewcommand{\ABNTEXsubsubsectionfont}{\rmfamily\itshape}
\renewcommand{\ABNTEXsubsubsectionfontsize}{\normalsize}

% ---------------------------------------------------------------- Capa e folha de rosto
% Versões próprias (NBR 14724): instituição, autores, título, subtítulo, local e ano.
\newcommand{\subtitulodoc}{}
\newcommand{\definirsubtitulo}[1]{\renewcommand{\subtitulodoc}{#1}}
\makeatletter
\newcommand{\capagrupo}{%
  \begin{capa}
    \centering
    {\bfseries\imprimirinstituicao\par}
    \vfill
    {\bfseries\large\MakeUppercase{\imprimirtitulo}\par}
    \ifx\subtitulodoc\@empty\else
      \vspace{0.4cm}{\large\subtitulodoc\par}\fi
    \vfill
    {\imprimirlocal\par}
    {\imprimirdata\par}
    \vspace*{0.5cm}
  \end{capa}}
\newcommand{\folhaderostogrupo}{%
  \begin{folhaderosto}
    \centering
    \vspace*{\fill}
    {\bfseries\large\MakeUppercase{\imprimirtitulo}\par}
    \ifx\subtitulodoc\@empty\else
      \vspace{0.4cm}{\large\subtitulodoc\par}\fi
    \vspace*{\fill}
    \hspace{.45\textwidth}%
    \begin{minipage}{.5\textwidth}
      \SingleSpacing\small
      \imprimirpreambulo
      \par\vspace{0.6cm}
      % nomes na mesma linha, separados por vírgula (\\ de \autor vira ",")
      \textbf{Discentes:} {\def\\{,}\imprimirautor}.\par
      \vspace{0.6cm}
      \textbf{\imprimirorientadorRotulo}~\imprimirorientador
    \end{minipage}\par
    \vspace*{\fill}
    {\imprimirlocal\par}
    {\imprimirdata\par}
    \vspace*{0.5cm}
  \end{folhaderosto}}
\makeatother
\newcommand{\listafiguras}{%
  \pdfbookmark[0]{\listfigurename}{lof}%
  \listoffigures*%
  \cleardoublepage}
\newcommand{\listatabelas}{%
  \pdfbookmark[0]{\listtablename}{lot}%
  \listoftables*%
  \cleardoublepage}
% Página pré-textual com título centralizado e sem número de seção
\newcommand{\paginapretextual}[1]{%
  \pdfbookmark[0]{#1}{\detokenize{#1}}%
  \begin{center}{\ABNTEXchapterfont\MakeUppercase{#1}}\end{center}\vspace{0.8cm}}
\newcommand{\listaquadros}{%
  \pdfbookmark[0]{\listofquadrosname}{loq}%
  \listofquadros*%
  \cleardoublepage}
% Legenda de quadro fora de float (para quadros longos com xltabular)
\newfixedcaption{\legendaquadro}{quadro}
\newcommand{\sumario}{%
  \pdfbookmark[0]{\contentsname}{toc}%
  \tableofcontents*%
  \cleardoublepage}

% Dados da capa e da folha de rosto (comuns aos três documentos).
% Autores na ordem do enunciado (Grupo VI). Para completar os nomes, edite apenas \autor abaixo.
\instituicao{%
  UNIVERSIDADE FEDERAL DO SUL E SUDESTE DO PARÁ -- UNIFESSPA
  \par
  FACULDADE DE SISTEMAS DE INFORMAÇÃO -- FACSI
  \par
  CURSO DE SISTEMAS DE INFORMAÇÃO}
\autor{%
  Youri\\
  José Miguel\\
  Mateus Leitão\\
  Michelangelo\\
  Ângelo}
\orientador[Professor:]{Gleison de Oliveira Medeiros}
\local{Marabá -- PA}
\data{2026}
\newcommand{\naturezatrabalho}{Trabalho~1 da disciplina Sistemas Distribuídos -- Engenharia de Requisitos e Desenvolvimento de um Agregador Concorrente (Tema~6: Agregador de Preços de Eletrônicos), apresentado ao Curso de Sistemas de Informação da Faculdade de Sistemas de Informação da Universidade Federal do Sul e Sudeste do Pará. Grupo~VI.}

\titulo{Agregador de Preços de Eletrônicos}
\definirsubtitulo{Modelagem do Sistema}
\preambulo{\naturezatrabalho}

\begin{document}
\selectlanguage{brazil}
\frenchspacing

% ================================================================= PRÉ-TEXTUAIS
\pretextual
\capagrupo
\folhaderostogrupo
\listafiguras
\listaquadros
\sumario

% ================================================================= TEXTUAIS
\textual

\chapter{Introdução}

Este documento reúne a modelagem do \textbf{Agregador de Preços de Eletrônicos} (Trabalho~1 de Sistemas Distribuídos, Tema~6, Grupo~VI), conforme a seção~5.4 do enunciado \cite{enunciado2026}: o \textbf{diagrama de casos de uso UML}, o \textbf{diagrama de arquitetura da aplicação} e o \textbf{diagrama de sequência de uma operação principal} (a pesquisa de um produto). Como pede o enunciado, os diagramas representam os componentes -- \textbf{cliente}, \textbf{agregador} e \textbf{serviços independentes} (seção~6 do enunciado) -- e os fluxos de comunicação entre eles. Os requisitos que esses diagramas atendem estão no Documento de Requisitos, e a implementação e os experimentos, no Relatório Técnico.

O diagrama de casos de uso segue a notação UML \cite{sommerville2015}. Os diagramas de arquitetura e de sequência também são mantidos como fontes Mermaid no repositório (\cod{docs/diagramas/*.mmd}); as figuras deste documento foram redesenhadas para impressão, com o mesmo conteúdo.

O sistema é composto por processos independentes que se comunicam apenas por \textbf{sockets TCP}: clientes, um agregador e quatro lojas simuladas. Cada mensagem é um objeto JSON em UTF-8, em uma linha terminada por \cod{\textbackslash n}, com uma requisição e uma resposta por conexão. A especificação completa do protocolo está em \cod{docs/03\_protocolo.md}.

% -----------------------------------------------------------------
\chapter{Diagrama de casos de uso}

A \autoref{fig:casos-uso} apresenta o diagrama UML de casos de uso. Os atores (bonecos) ficam fora da fronteira do sistema e os casos de uso são elipses. As associações entre atores e casos de uso são linhas contínuas; as relações «include» e «extend» são linhas tracejadas com seta aberta. O ator secundário \textit{Serviço de loja} representa as quatro lojas.

\begin{figure}[htb]
  \centering
  \caption{Diagrama de casos de uso}\label{fig:casos-uso}
  \diagrama{casos_de_uso}
  \fonteautores
\end{figure}

O \autoref{qua:casos-uso} resume cada elemento do diagrama.

\begin{quadro}[!htbp]
  \centering\small
  \caption{Casos de uso do sistema}\label{qua:casos-uso}
  \begin{tabularx}{\textwidth}{|L{4.0cm}|L{3.4cm}|Y|}
    \hline
    \textbf{Caso de uso} & \textbf{Ator} & \textbf{Descrição resumida} \\ \hline
    UC01 -- Pesquisar produto & Usuário comprador; Serviço de loja & Pesquisa um eletrônico nas quatro lojas e apresenta as ofertas consolidadas. \\ \hline
    UC02 -- Comparar sequencial × concorrente & Usuário comprador & Executa o UC01 nos dois modos e mostra o speedup («include» UC01). \\ \hline
    UC03 -- Listar catálogo & Usuário comprador; Serviço de loja & Mostra os produtos disponíveis. \\ \hline
    UC04 -- Executar experimentos & Pesquisador & Roda os cinco experimentos e gera tabelas e gráficos («include» UC01). \\ \hline
    UC05 -- Usar a interface web & Usuário comprador; Pesquisador & Pesquisa, compara, derruba e sobe lojas, faz teste de carga e vê os experimentos pelo navegador («include» UC01 e UC02). \\ \hline
    Ordenar e identificar melhor oferta & -- & Incluído no UC01: ordena por preço total ou custo-benefício. \\ \hline
    Registrar tempo de execução & -- & Incluído no UC01: grava o tempo em \cod{log\_pesquisas.csv}. \\ \hline
    Tratar falha de loja & -- & Estende o UC01 quando uma loja está fora do ar, cai ou estoura o timeout. \\ \hline
  \end{tabularx}
  \fonteautores
\end{quadro}

% -----------------------------------------------------------------
\chapter{Diagrama de arquitetura}

A \autoref{fig:arquitetura} mostra a arquitetura em três camadas de processos: clientes $\rightarrow$ agregador (porta 8000) $\rightarrow$ quatro lojas (portas 9001--9004). Toda a comunicação entre eles usa TCP com mensagens JSON delimitadas por quebra de linha. O navegador fala HTTP apenas com o gateway (porta 8080), que é mais um cliente TCP do agregador.

\begin{figure}[!htbp]
  \centering
  \caption{Diagrama de arquitetura}\label{fig:arquitetura}
  \diagrama{arquitetura}
  \fonteautores
\end{figure}

\section{Componentes}

O \autoref{qua:componentes} descreve cada componente, o processo em que executa, a porta e a concorrência interna. Todos os módulos Python ficam em \cod{backend/}.

\begin{quadro}[!htbp]
  \centering\small
  \caption{Componentes da arquitetura}\label{qua:componentes}
  \begin{tabularx}{\textwidth}{|L{3.6cm}|L{2.0cm}|L{2.2cm}|Y|}
    \hline
    \textbf{Componente} & \textbf{Processo} & \textbf{Porta} & \textbf{Concorrência interna} \\ \hline
    Cliente (\cod{cliente.py}) & um por usuário & efêmera & -- (com \cod{--clientes N} ou em \cod{experimentos.py}, N clientes em threads) \\ \hline
    Gateway HTTP (\cod{gateway\_http.py}) & 1 & 8080 (HTTP) & 1 thread por requisição HTTP; é um cliente TCP do agregador (no teste de carga, N clientes em threads) \\ \hline
    Interface web (\cod{frontend/}, Vite + React) & navegador & 5173 em desenvolvimento & -- (fala só HTTP com o gateway) \\ \hline
    Agregador (\cod{agregador.py}) & 1 & 8000 & Thread de aceite + 1 thread por cliente; no modo concorrente, +1 thread por loja \\ \hline
    Lojas (\cod{loja\_server.py} × 4) & 4 independentes & 9001--9004 & Thread de aceite + 1 thread por conexão \\ \hline
    Catálogos (\cod{dados/lojas/*.json}) & -- & -- & Lidos uma vez quando a loja inicia \\ \hline
  \end{tabularx}
  \fonteautores
\end{quadro}

\section{Módulos de apoio}

Os módulos do \autoref{qua:modulos} não são processos próprios: são usados pelos componentes acima.

\begin{quadro}[!htbp]
  \centering\small
  \caption{Módulos de apoio}\label{qua:modulos}
  \begin{tabularx}{\textwidth}{|L{3.8cm}|Y|}
    \hline
    \textbf{Módulo} & \textbf{Função} \\ \hline
    \cod{protocolo.py} & Envio e recebimento das mensagens (\textit{framing} JSON + \cod{\textbackslash n}, limite de 64~KB), operações, modos, critérios e códigos de status. \\ \hline
    \cod{servidor\_base.py} & Servidor TCP multithread genérico (\cod{ServidorTCPMultithread}), usado pelo agregador e pelas lojas. \\ \hline
    \cod{config.py} & Endereços, portas, latências e timeouts. \\ \hline
    \cod{iniciar\_servicos.py} & Sobe as cinco instâncias de servidor (quatro lojas + agregador) como processos e espera que respondam a \cod{ping}; também é usado pelo gateway. \\ \hline
  \end{tabularx}
  \fonteautores
\end{quadro}

\section{Fluxo de comunicação}

O cliente abre uma conexão TCP com o agregador e envia uma requisição. A thread do cliente no agregador escolhe a estratégia:

\begin{itemize}
  \item \textbf{sequencial}: um laço consulta loja por loja;
  \item \textbf{concorrente}: uma thread por loja, todas aguardadas com \cod{join()}.
\end{itemize}

Cada consulta a uma loja é uma conexão TCP própria, com timeout de 2,5~s. As respostas são consolidadas, o tempo é registrado no CSV (protegido por \cod{Lock}) e o resultado volta ao cliente na mesma conexão. Como cada pesquisa abre conexões novas, o protocolo não guarda estado entre pesquisas: uma loja que volta ao ar passa a ser consultada de novo sem reiniciar o agregador \cite{tanenbaum2007}.

\section{Interface web}

O navegador não abre sockets TCP, então a interface chama o gateway por HTTP/JSON (\cod{/api/...}). O gateway valida o pedido e usa a mesma função do \cod{cliente.py} (\cod{requisitar\_agregador}) para abrir uma conexão TCP com o agregador. Do agregador para dentro nada muda: ele não sabe se o pedido veio da linha de comando ou da interface. A API HTTP está no Apêndice~A de \cod{docs/03\_protocolo.md}.

% -----------------------------------------------------------------
\chapter{Diagrama de sequência}

A \autoref{fig:sequencia} mostra o caso de uso ``Pesquisar produto'' no modo concorrente, com uma loja fora do ar.

\begin{figure}[!htbp]
  \centering
  \caption{Diagrama de sequência -- pesquisa no modo concorrente, com uma loja fora do ar}\label{fig:sequencia}
  \diagrama{sequencia}
  \fonteautores
\end{figure}

O bloco \textbf{par} mostra as quatro consultas executadas \textbf{ao mesmo tempo}, em threads diferentes do agregador. A thread~4 ilustra o tratamento de falha: a conexão com a MagaluTechSim é recusada e o agregador registra \cod{erro\_comunicacao} para essa loja, sem atrasar as demais. Depois do \cod{join()}, o agregador ordena as ofertas pelo preço total, aponta o menor preço (PichauSim, R\$~1.864,00) e grava a pesquisa em \cod{log\_pesquisas.csv}. O tempo total corresponde ao da consulta mais lenta, e não à soma.

No \textbf{modo sequencial}, o diagrama seria igual, mas sem o bloco \textbf{par}: as mensagens 3 a 9 aconteceriam uma após a outra, e cada consulta só começaria depois da resposta da anterior. Por isso o tempo total do modo sequencial se aproxima da soma das latências das lojas, como mostram os experimentos do Relatório Técnico.

% ================================================================= PÓS-TEXTUAIS
\postextual
\bibliography{referencias}

\end{document}
